
\section{The proof-carrying module calculus}
\label{sec:modules}

\subsection{The basic record}

Fix a free-variable scope \(n\).  A module has a finite list of requirements
and one conclusion.

\begin{definition}[PCFM module]
A PCFM module in scope \(n\) is a triple
\[
   M=(\mathsf{req}_M,\mathsf{out}_M,\pi_M)
\]
where
\[
  \mathsf{req}_M:\mathsf{List}(\Fm\,\PTm\,n),\qquad
  \mathsf{out}_M:\Fm\,\PTm\,n,
\]
and
\[
  \pi_M:
  \mathsf{req}_M\vdash_{\PRAAx}\mathsf{out}_M.
\]
\end{definition}

Requirements are not divided by the kernel into ``ordinary assumptions'' and
``residual obligations''.  That is metadata used by the front end.  Logically,
both are formulas in the context.  This is useful: the same simultaneous-cut
operation discharges either kind.

\subsection{Specialization and weakening}

\begin{theorem}[Specialization]
Let \(M\) be a module in scope \(n\), and let
\(\sigma:\mathsf{Fin}\,n\to\PTm\,m\) be a simultaneous term substitution.
Then there is a module \(M[\sigma]\) in scope \(m\) with
\[
  \mathsf{req}_{M[\sigma]}
   = \{\psi[\sigma]:\psi\in\mathsf{req}_M\},
  \qquad
  \mathsf{out}_{M[\sigma]}=\mathsf{out}_M[\sigma].
\]
\end{theorem}

\begin{proof}
Apply \code{Prf.subst} to \(\pi_M\).  The theorem already proves
capture-avoiding simultaneous substitution through every connective and
quantifier, so no module-specific binder argument is needed.
\end{proof}

\begin{theorem}[Weakening]
If every formula in \(\mathsf{req}_M\) belongs to a context \(\Delta\), then
\[
   \Delta\vdash_{\PRAAx}\mathsf{out}_M.
\]
\end{theorem}

\begin{proof}
Apply \code{Prf.weaken} to \(\pi_M\).  Contexts are lists, but the structural
theorem uses list inclusion and hence ignores order and duplication.
\end{proof}

These two operations are already checked in \code{ConciseIPRA.Deduction}.  The
PCFM implementation should re-export them with names appropriate to theorem
assembly rather than reproving them.

\subsection{Simultaneous discharge}

\begin{theorem}[Module discharge]
Let \(M\) have requirements \(\Gamma\).  Suppose a context \(\Delta\) derives
every \(\psi\in\Gamma\).  Then
\[
   \Delta\vdash_{\PRAAx}\mathsf{out}_M.
\]
\end{theorem}

\begin{proof}
This is exactly \code{Prf.trans}.  Its proof repeatedly discharges context
formulas by implication introduction and restores their supplied derivations by
modus ponens.  It is a simultaneous cut theorem internal to the checked
natural-deduction metatheory.
\end{proof}

The theorem is the engine of residualization.  If \(M\) proves
\[
  B_1,\ldots,B_r,A_1,\ldots,A_s\vdash C
\]
and modules \(N_i\) prove the \(B_i\) and \(K_j\) prove the \(A_j\) from a
common context \(\Delta\), then one call to \code{Prf.trans} produces
\(\Delta\vdash C\).

\subsection{Unary contracts and categorical composition}

For readability, many nodes use one precondition.

\begin{definition}[Contract]
A contract \(A\Rightarrow B\) is a derivation
\[
       [A]\vdash_{\PRAAx}B.
\]
\end{definition}

\begin{theorem}[Identity and composition]
Contracts form a category:
\[
  \mathsf{id}_A:[A]\vdash A,
\]
and from \(f:[A]\vdash B\) and \(g:[B]\vdash C\) one obtains
\[
  g\circ f:[A]\vdash C.
\]
Composition is associative and identities are neutral as derivations up to the
propositional equality induced by their common conclusion.
\end{theorem}

\begin{proof}
Identity is the hypothesis rule.  For composition, apply
\code{Prf.trans} to \(g\), replacing its only hypothesis \(B\) by \(f\).
Associativity is not required as an equality of proof syntax: either
parenthesization has type \([A]\vdash C\).  If a definitional equality of
compiler outputs is desired, it can be established separately; mathematical
soundness needs only derivability.
\end{proof}

Multi-input contracts may either use a list context or package the inputs into
a conjunction.  Lists are preferable for residual producer obligations
because they can be discharged independently.  Conjunctions are useful when
the formula itself is to be passed as a single node output.

\subsection{Exact equations as rewrite modules}

An equation proof
\[
   \Gamma\vdash s=t
\]
may be transported through any term or formula using the existing equality
rules.  PCFM should expose specialized constructors:

\[
\begin{array}{ll}
\mathsf{congrApp} &: (t_i=u_i)_i\mapsto f(\vec t)=f(\vec u),\\
\mathsf{rewriteTerm} &: s=t,\ h(s)=v\mapsto h(t)=v,\\
\mathsf{rewriteFormula} &: s=t,\ \Phi(s)\mapsto\Phi(t),\\
\mathsf{chainEq} &: s=t,\ t=u\mapsto s=u.
\end{array}
\]

The current derivation library already supplies congruence for application,
addition, multiplication, predecessor, truncated subtraction, and many
derived symbols.  The module layer should depend on those theorems rather than
reimplementing algebra.

\subsection{Residual contexts and provenance}

A requirement should carry host-level metadata:
\[
   (\mathsf{name},\mathsf{kind},\mathsf{source},\mathsf{formula}).
\]
The metadata is not part of the trusted derivation.  Its purpose is audit and
automation.  Recommended kinds are:

\begin{center}
\begin{tabular}{ll}
\toprule
Kind & Meaning\\
\midrule
\code{local} & ordinary theorem input or previously proved fact\\
\code{resource} & bound, modulus, clock, finite cover, or precision demand\\
\code{separation} & apartness, nonzero denominator, isolation, orientation\\
\code{certificate} & validity of supplied finite evidence\\
\code{external} & semantic input not yet internalized\\
\code{bridge} & formula/code correspondence needed for backend transport\\
\bottomrule
\end{tabular}
\end{center}

The compiled theorem stores only the formulas.  A separate ledger records
which module discharged each requirement.  This prevents a proof from silently
changing status when an external assumption is renamed as a definition.

\subsection{The local soundness theorem is enough}

A node implementation need not establish a global semantic theorem about an
entire graph.  It establishes its local contract in the object calculus.
Global soundness follows by repeated discharge.

\begin{theorem}[Acyclic assembly]
Let \(G\) be a finite acyclic dependency graph.  Associate to each node \(v\)
a module whose requirements are formulas produced by predecessor nodes or
listed residuals.  If every predecessor edge supplies a derivation of the
corresponding requirement, then the sink formula is derivable from precisely
the undischarged residual requirements.
\end{theorem}

\begin{proof}
Choose a topological order.  Induct through that finite list at the host level.
At node \(v\), predecessor outputs have already been compiled.  Apply module
discharge to the local proof.  Cache the resulting derivation as the output of
\(v\).  Residual formulas are left in the common context and are available by
the hypothesis rule.  The sink derivation has no other assumptions because
each nonresidual requirement was replaced by a predecessor derivation.
\end{proof}

This is a metatheorem about the Lean assembler.  The trusted output is not the
metatheorem but the generated \code{Prf PRAAx} value.  Consequently even a bug
in graph scheduling cannot certify a false theorem: it can at worst fail to
construct a term of the required type.

\subsection{Why this is a substantial simplification}

A monolithic compiler theorem must formalize syntax, typing, evaluation,
substitution, environments, traces, and graph soundness before it produces one
application proof.  PCFM reuses the existing object syntax and deduction
metatheory.  Its structural core is almost entirely aliases for already
checked lemmas:

\[
\begin{array}{c|c}
\text{PCFM operation} & \text{existing checked theorem}\\
\hline
\text{specialize} & \code{Prf.subst}\\
\text{weaken} & \code{Prf.weaken}\\
\text{discharge/cut} & \code{Prf.trans}\\
\text{open a binder} & \code{Prf.allOpen}\\
\text{instantiate} & \code{Prf.allE},\ \code{IPRA.inst}\\
\text{quantifier-free induction} & \code{PRA.qfInd},\ \code{IPRA.ind}\\
\text{arithmetic recursion} & \code{recDef},\ \code{rec_zero},\ \code{rec_succ}
\end{array}
\]

The genuinely new work is therefore concentrated in Boolean reflection,
generic trace packaging, viability, stable observation, and public bridge
wrappers.
